\subsection{有限 Hardy 证书、有理主弧 Gram 核与 phase Fourier--Mellin 双轴刚性}\label{subsec:conclusion-phase-gram-mellin-hardy-radial-rigidity}
沿用 circle-dimension phase-gate 一章关于有限模态 Hardy 投影、有理主弧 Gram 核、双底去周期化、尺度 Mellin 叠加完成化、Mellin--Plancherel 幺正切片与连续横截方向唯一性的记号。前文已经分别证明：有限模态相位投影可由显式有理核完全实现；主弧 Gram 二次型的谱贡献可被严格局域化到低频 Ramanujan 相关带；一次尺度上的谱降会通过尺度自举传播为全尺度幂次衰减，且统一谱隙的失败必先出现在有限检查集；两条对数不可公度尺度已足以摧毁单尺度强制的严格虚周期；尺度 Mellin 叠加在保持函数方程与临界线实值性的同时允许连续底混合；最后，unweighted 的等距 Mellin 读出只能落在临界线，而纯相位制度中唯一允许的连续横截方向只能是径向寄存器。把这些接口进一步并读后，可压出如下 phase Fourier--Mellin 闭链。

\begin{proposition}[有限模态 Hardy 证书的有理再生性]\label{prop:conclusion-finite-mode-hardy-rational-reproducibility}
\leanverified{paper\_conclusion\_finite\_mode\_hardy\_rational\_reproducibility}
令
\[
w(x):=\frac{1+\mathrm{i}x}{1-\mathrm{i}x},
\]
并令
\[
\Pi_N:L^2(\RR,\dd\mu)\to \mathrm{Span}\{w(\cdot)^k:0\le k\le N-1\}
\]
为前 \(N\) 个非负模态的 Hardy 正交投影。则 \(\Pi_N\) 由显式有理核
\[
K_N(x,y)
=
\sum_{k=0}^{N-1}w(x)^k\overline{w(y)}^{\,k}
=
\frac{1-\bigl(w(x)\overline{w(y)}\bigr)^N}{1-w(x)\overline{w(y)}}
\]
唯一表示；并且当 \(x\neq y\) 时有分母闭式化简
\[
\frac{1}{1-w(x)\overline{w(y)}}
=
\frac12+\frac{1+xy}{2\mathrm{i}(y-x)}.
\]
因此，每一个有限模态 Hardy 证书都同时满足：
\begin{enumerate}
  \item \(\operatorname{rank}\Pi_N=N\)；
  \item 投影核完全由有理函数给出；
  \item 全部非平凡奇性都被压缩到单一 Cauchy 型分母上，而不分散到高模态尾部。
\end{enumerate}
\end{proposition}
\begin{proof}
命题 \ref{prop:unit-circle-phase-hardy-kernel-rational} 已给出核公式
\eqref{eq:unit-circle-phase-hardy-kernel}、分母化简
\eqref{eq:unit-circle-phase-hardy-kernel-denominator} 以及积分算子表示
\[
(\Pi_N f)(x)=\int_{\RR}K_N(x,y)f(y)\,\dd\mu(y).
\]
由于 \(\Pi_N\) 是到 \(N\)-维模态子空间的正交投影，故 \(\operatorname{rank}\Pi_N=N\)。核表示中的有限几何级数与分母闭式说明其全部信息都落在显式有理表达式中，而非平凡奇性又只由 \((y-x)^{-1}\) 这一单一 Cauchy 通道承载。证毕。
\end{proof}

\begin{corollary}[有限模态 Hardy 核的平坦对角与精确有限秩]\label{cor:conclusion-finite-mode-hardy-flat-diagonal-exact-rank}
\leanverified{paper\_conclusion\_finite\_mode\_hardy\_flat\_diagonal\_exact\_rank}
沿用命题 \ref{prop:conclusion-finite-mode-hardy-rational-reproducibility} 的记号。则对每个 \(N\ge 1\) 与每个 \(x\in\RR\) 都有
\[
K_N(x,x)=N,
\]
并且
\[
\dim_{\CC}\operatorname{Ran}\Pi_N=N.
\]
因此，有限模态 Hardy 证书不仅由显式有理核承载，而且其 Christoffel 对角在整个实轴上完全平坦，像空间复维数也被该同一参数 \(N\) 精确锁定。
\end{corollary}
\begin{proof}
由命题 \ref{prop:unit-circle-phase-hardy-kernel-rational} 的对角值公式，
\[
K_N(x,x)=\sum_{k=0}^{N-1}|w(x)|^{2k}=N,
\]
因为 \(w(x)\in\TT\) 从而 \(|w(x)|=1\)。另一方面，\(\Pi_N\) 的像空间按定义就是
\[
\operatorname{span}\{w(\cdot)^k:0\le k\le N-1\},
\]
故复维数恰为 \(N\)。证毕。
\end{proof}

\begin{theorem}[主弧 Gram 核的 Ramanujan 局域化]\label{thm:conclusion-majorarc-gram-ramanujan-localization}
记
\[
Q_M:=\left\{\frac aq:1\le q\le M,\ (a,q)=1\right\}\subset \QQ/\ZZ,
\qquad
\ell^2(Q_M)\cong \bigoplus_{q\le M}\ell^2\bigl((\ZZ/q\ZZ)^\times\bigr).
\]
取一个实偶、非负、快速衰减核 \(\widehat\psi\)，令 \(\psi_L\) 为其 \(L\)-尺度周期化核，并记相应主弧 Gram 算子为 \(K_{M,L}\)。则对任意 \(g\in\ell^2(Q_M)\) 都有严格恒等式
\[
\langle g,K_{M,L}g\rangle
=
\sum_{n\in\ZZ}\widehat\psi(n/L)
\left|
\sum_{q\le M}\ \sum_{(a,q)=1}g(a/q)e^{2\pi\mathrm{i}an/q}
\right|^2.
\]
因此，\(\|K_{M,L}\|_{\mathrm{op}}\) 的主要贡献被完全压缩到一条低频 Ramanujan 相关带中；尤其当 \(g\) 在每个分母层上近常数时，内层和退化为 Ramanujan 和 \(c_q(n)\)，从而整个谱范数问题下降为有限分母层的低频相关审计。
\end{theorem}
\begin{proof}
这正是命题 \ref{prop:conclusion66-gram-ramanujan-decomposition} 的幅度平方恒等式。其右端按频率 \(n\) 精确分解为非负平方振幅，因此谱范数主贡献不可能隐藏在某种连续尾部，而必然由有限分母层与有限低频带之间的相关结构承载。若 \(g\) 在每个分母层上为常数，则内层和直接退化为 Ramanujan 和 \(c_q(n)\)，故问题完全下降为算术相关审计。证毕。
\end{proof}
\leanverified{paper\_conclusion\_majorarc\_gram\_ramanujan\_localization}

\begin{corollary}[分母层常值测试向量的 Ramanujan 相关压缩]\label{cor:conclusion-majorarc-gram-denominator-constant-ramanujan-correlation}
在定理 \ref{thm:conclusion-majorarc-gram-ramanujan-localization} 的设定下，若 \(g\in \ell^2(Q_M)\) 仅依赖于分母，即
\[
g(a/q)=g_q
\qquad
(1\le q\le M,\ (a,q)=1),
\]
则有严格恒等式
\[
\langle g,K_{M,L}g\rangle
=
\sum_{n\in\ZZ}\widehat\psi(n/L)
\left|
\sum_{q\le M} g_q\,c_q(n)
\right|^2,
\]
其中
\[
c_q(n):=\sum_{(a,q)=1}e^{2\pi\mathrm{i}an/q}
\]
为 Ramanujan 和。故在分母层常值口径下，主弧 Gram 二次型的主贡献被精确压缩到有限分母层的低频 Ramanujan 相关之中。
\end{corollary}
\begin{proof}
将
\(
g(a/q)=g_q
\)
直接代入定理 \ref{thm:conclusion-majorarc-gram-ramanujan-localization} 的幅度平方恒等式，可得
\[
\sum_{q\le M}\ \sum_{(a,q)=1}g(a/q)e^{2\pi\mathrm{i}an/q}
=
\sum_{q\le M} g_q \sum_{(a,q)=1}e^{2\pi\mathrm{i}an/q}
=
\sum_{q\le M} g_q\,c_q(n).
\]
代回即得结论。证毕。
\end{proof}

\begin{theorem}[有限算术证书对全尺度谱隙的完备控制]\label{thm:conclusion-finite-arithmetic-certificates-control-fullscale-gap}
\leanverified{paper\_conclusion\_finite\_arithmetic\_certificates\_control\_fullscale\_gap}
在定理 \ref{thm:conclusion-majorarc-gram-ramanujan-localization} 的框架下，取标度关系 \(L\asymp M^2\)。若存在某个尺度 \(M_0\) 使
\[
\|K_{M_0,L_0}\|_{\mathrm{op}}\le M_0^{-\eta}
\qquad
(\eta>0,\ L_0\asymp M_0^2),
\]
则存在 \(\eta'>0\) 使对所有充分大的 \(M\) 及 \(L\asymp M^2\) 都有
\[
\|K_{M,L}\|_{\mathrm{op}}\ll M^{-\eta'}.
\]
反之，若全尺度统一谱隙失败，则必在一个有限可枚举检查集
\[
\mathcal C\subset (0,2\pi)\times \NN
\]
中出现可复核反例。因而主弧 Gram 核的全尺度谱隙不是无限验证问题，而是一个有限算术证书问题。
\end{theorem}
\begin{proof}
定理 \ref{thm:conclusion67-scale-bootstrap} 已表明：在 \(L\asymp M^2\) 标度下，一次尺度上的谱降会自举为全尺度的幂次衰减。另一方面，定理 \ref{thm:conclusion68-finite-verification} 已把统一扭曲谱隙等价归约为有限检查集上的核验。将这两条与定理 \ref{thm:conclusion-majorarc-gram-ramanujan-localization} 的 Ramanujan 局域化恒等式并置，便得到所述结论：一次多项式节省会传播到全尺度，而任何统一谱隙的失败都不可能藏在无限远尺度中，而必先在某个有限检查对上显影。证毕。
\end{proof}

\begin{theorem}[两尺度去周期化的最小性]\label{thm:conclusion-two-incommensurable-scales-minimal-deperiodization}
设 \(L_1,L_2>1\)，并令
\[
f(s):=g(L_1^{-s})+h(L_2^{-s}),
\]
其中 \(g,h\) 不具有任何非平凡单位根乘法不变性。若
\[
\frac{\log L_1}{\log L_2}\notin\QQ,
\]
则 \(f\) 不存在任何非零纯虚周期。反之，单一尺度族
\[
s\longmapsto g(L^{-s})
\]
必携带严格虚周期梳齿。
因此，两条不可公度尺度正是摧毁严格虚周期的最小机制。
\end{theorem}
\begin{proof}
正向断言即定理 \ref{thm:xi-two-base-deperiodization-no-imag-period}。反向断言则由命题 \ref{prop:single-scale-dirichlet-imag-period} 给出：单尺度 Dirichlet 化总是具有严格虚周期
\[
T_L=\frac{2\pi\mathrm{i}}{\log L}.
\]
因此，少于两条不可公度尺度时，竖向周期性无法被摧毁；而一旦引入两条对数不可公度尺度，则任何非零纯虚周期都被禁止。证毕。
\end{proof}

\begin{theorem}[对称去周期化与函数方程保持的可积完备性]\label{thm:conclusion-scale-mellin-deperiodization-preserves-selfdual-reality}
在定理 \ref{thm:conclusion-two-incommensurable-scales-minimal-deperiodization} 的双尺度极小去周期机制之上，令 \(w\) 为实值权函数，并定义尺度 Mellin 叠加完成化
\[
\mathcal X_{\Delta,w}(s)
:=
\int_1^\infty w(\log L)\,\Xi_{\Delta,L}(s)\,\frac{\dd L}{L}.
\]
则：
\begin{enumerate}
  \item \(\mathcal X_{\Delta,w}\) 保持与单尺度完成化相同的自对偶函数方程
  \[
  \mathcal X_{\Delta,w}(s)=\mathcal X_{\Delta,w}(1-s);
  \]
  \item 在临界线 \(\Re(s)=\frac12\) 上保持实值性
  \[
  \mathcal X_{\Delta,w}\!\left(\frac12+\mathrm{i}t\right)\in\RR;
  \]
  \item 因而，双尺度去周期化与连续底混合并不破坏“自对偶--临界线实值”制度。
\end{enumerate}
\end{theorem}
\begin{proof}
命题 \ref{prop:xi-scale-mellin-fe-real} 已给出尺度 Mellin 叠加完成化的函数方程与临界线实值性。故连续底混合并不以牺牲自对偶结构为代价，而是在保持该对称制度的前提下引入去周期化所需的连续尺度自由度。证毕。
\end{proof}

\begin{theorem}[临界线的等距唯一性]\label{thm:conclusion-critical-line-unique-unweighted-isometric-slice}
定义 Mellin 变换
\[
(\mathcal M f)(s):=\int_0^\infty f(x)\,x^{s-\frac12}\,\dd^\times x,
\qquad
\dd^\times x=\frac{\dd x}{x},
\]
并记
\[
(\mathcal M_\sigma f)(t):=(\mathcal M f)(\sigma+\mathrm{i}t).
\]
则：
\begin{enumerate}
  \item 切片
  \[
  \mathcal M_{1/2}:L^2_\times\to L^2\!\left(\RR,\frac{\dd t}{2\pi}\right)
  \]
  是幺正同构；
  \item 对任意 \(\sigma\in\RR\)，都有严格权重扭曲恒等式
  \[
  \|\mathcal M_\sigma f\|_{L^2(\dd t/2\pi)}
  =
  \|x^{\sigma-\frac12}f\|_{L^2_\times};
  \]
  \item 因而，唯一不需要额外权重寄存器的等距切片恰是
  \[
  \Re(s)=\frac12.
  \]
\end{enumerate}
\end{theorem}
\begin{proof}
这正是定理 \ref{thm:mellin-unitary-slice-unique} 的内容。第一条给出 \(\sigma=\frac12\) 时的幺正性；第二条说明所有其它切片都必须引入额外权重 \(x^{\sigma-\frac12}\)。因此，若坚持“不外置任何额外权重记录轴”，则唯一允许的等距读出切片只能是临界线 \(\Re(s)=\frac12\)。证毕。
\end{proof}

\begin{theorem}[连续横截方向的唯一性]\label{thm:conclusion-radial-axis-unique-continuous-transverse-direction}
设纯相位读出为
\[
u\longmapsto \operatorname{proj}_{\mathsf{Phase}}(u).
\]
若允许外置一个连续记录轴 \(r\) 使
\[
u\longmapsto \bigl(\operatorname{proj}_{\mathsf{Phase}}(u),r(u)\bigr)
\]
成为单射，则任意这样的连续记录轴都必经由 \(|u|\) 因子化；等价地，它与径向寄存器
\[
\varrho(u):=\log|u|
\]
仅差单调重参数化。
因此，在纯相位制度中，相位轴之外唯一允许的连续横截方向就是径向轴。
\end{theorem}
\begin{proof}
这正是推论 \ref{cor:xi-unique-continuous-transverse-register} 的内容：任何使纯相位读出单射化的连续记录轴都必因子化经由 \(|u|\)，从而与 \(\varrho(u)=\log|u|\) 只差一个连续单调重参数化。证毕。
\end{proof}

\begin{corollary}[相位--径向二轴的范畴刚性]\label{cor:conclusion-phase-radial-two-axis-categorical-rigidity}
由定理 \ref{thm:conclusion-critical-line-unique-unweighted-isometric-slice} 与定理 \ref{thm:conclusion-radial-axis-unique-continuous-transverse-direction} 可知：
\begin{enumerate}
  \item 唯一的内生等距读出方向是临界线切片 \(\Re(s)=\frac12\)；
  \item 唯一的可外置连续横截方向是径向寄存器 \(\varrho(u)=\log|u|\)；
  \item 因而，任何同时保持自对偶、临界线实值与扩展保存的连续制度，都必然分裂为
  \[
  \boxed{
  \text{临界线上的幺正相位读出}
  \;\oplus\;
  \text{径向记录轴}
  }.
  \]
\end{enumerate}
\end{corollary}
\begin{proof}
第一点由定理 \ref{thm:conclusion-critical-line-unique-unweighted-isometric-slice}；第二点由定理 \ref{thm:conclusion-radial-axis-unique-continuous-transverse-direction}；第三点是二者的直接合并。证毕。
\end{proof}

\begin{theorem}[phase Fourier--Mellin 闭链]\label{thm:conclusion-phase-fourier-mellin-closed-chain}
第 \(12.26\) 至 \(12.41\) 节所蕴含的有限与连续结构可被压缩为一条严格闭合链：
\[
\boxed{
\text{有限模态 Hardy 投影}
\Longrightarrow
\text{有理主弧 Gram 核}
\Longrightarrow
\text{单尺度谱降的全尺度传播}
\Longrightarrow
\text{有限算术证书的完备控制}
\Longrightarrow
\text{最小两尺度去周期化}
\Longrightarrow
\text{连续底 Mellin 完成化}
\Longrightarrow
\text{临界线等距唯一性}
\Longrightarrow
\text{径向横截唯一性}.
}
\]
更准确地说：
\begin{enumerate}
  \item 有限模态截断由显式有理核 \(K_N(x,y)\) 完全实现，故连续对象已在有限秩层被代数化；
  \item 主弧 Gram 核把谱范数问题严格局域化为有限分母层的 Ramanujan 相关和；
  \item 一次尺度上的谱降经尺度自举立刻传播到全尺度，而统一谱隙又可由有限检查集完全判定；
  \item 双底不可公度叠加是摧毁严格虚周期的最小机制；
  \item 尺度 Mellin 叠加完成化保持函数方程与临界线实值性；
  \item Mellin--Plancherel 唯一性把“无外置权重的等距读出”唯一钉在临界线；
  \item 扩展保存制度又把唯一连续横截方向钉在径向轴上。
\end{enumerate}
因此，这一段落中的有限核、Ramanujan 相关、Mellin 完成与记录轴并非几套平行字典，而是已经内生闭合成一条新的刚性链：
\[
\boxed{
\text{算术有限证书}
\Longrightarrow
\text{连续核谱}
\Longrightarrow
\text{最小去周期化}
\Longrightarrow
\text{临界线唯一性}
\Longrightarrow
\text{径向轴唯一性}.
}
\]
\end{theorem}
\begin{proof}
逐条应用命题 \ref{prop:conclusion-finite-mode-hardy-rational-reproducibility}、定理 \ref{thm:conclusion-majorarc-gram-ramanujan-localization}、定理 \ref{thm:conclusion-finite-arithmetic-certificates-control-fullscale-gap}、定理 \ref{thm:conclusion-two-incommensurable-scales-minimal-deperiodization}、定理 \ref{thm:conclusion-scale-mellin-deperiodization-preserves-selfdual-reality}、定理 \ref{thm:conclusion-critical-line-unique-unweighted-isometric-slice}、定理 \ref{thm:conclusion-radial-axis-unique-continuous-transverse-direction} 及推论 \ref{cor:conclusion-phase-radial-two-axis-categorical-rigidity} 即得。证毕。
\end{proof}

\endinput
